\documentclass[11pt,a4paper]{article}
\usepackage[T1]{fontenc}
\usepackage[utf8]{inputenc}
\usepackage{lmodern}
\usepackage[margin=25mm]{geometry}
\usepackage{amsmath,amssymb,amsthm,mathtools}
\usepackage{enumitem}
\usepackage[hidelinks]{hyperref}
\numberwithin{equation}{section}
\newtheorem{theorem}{Theorem}[section]
\newtheorem{lemma}[theorem]{Lemma}
\newtheorem{proposition}[theorem]{Proposition}
\newtheorem{corollary}[theorem]{Corollary}
\theoremstyle{remark}\newtheorem{remark}[theorem]{Remark}
\newcommand{\R}{\mathbb R}
\newcommand{\cA}{\mathcal A}
\newcommand{\cD}{\mathcal D}
\newcommand{\avg}[1]{\left\langle #1\right\rangle}
\DeclareMathOperator{\inter}{int}
\DeclareMathOperator{\diver}{div}
\setlist{nosep,leftmargin=2.2em}
\title{The global attractor of the Wilhelm--Heinrich oscillator\\ is a convex $C^1$ graph over a closed Jordan domain}
\author{Anonymous (lane \texttt{planar}, working note)}
\date{3 October 2026}
\begin{document}
\maketitle

\noindent\textbf{Status.} Every step below is written out. The note proves the attractor-graph theorem
(Theorem~\ref{thm:A}) and its corollaries. It does \emph{not} prove uniqueness of the limit cycle; it reduces that
conjecture to a statement about a planar $C^1$ flow with a convex right-hand side. Imported results are listed in
\S\ref{sec:inputs}; nothing else is assumed. Floating-point checks of every numerical consequence are reported in
\S\ref{sec:numerics}; they are consistency tests and are not part of any proof.

\section{Setting, imported inputs and the theorem}\label{sec:setting}

Fix $\mu>0$, $k>0$ and put $A=1+\mu$, $\lambda=k/\mu$, $k_H=\mu(1+\mu)$. In logarithmic coordinates
$z=(U,p,y)\in\R^3$ the Wilhelm--Heinrich system is
\begin{equation}\label{eq:sys}
U'=p,\qquad p'=k-\mu p-y,\qquad y'=ke^{U}-y ,
\end{equation}
with vector field $f$, flow $\varphi_t$, $x:=ke^{U}$, equilibrium $E=(0,0,k)$, and positive chemical domain
$\cD=\{(U,p,y):y>0,\ p<\lambda\}$ (that is $X=x/k>0$, $Z=y/k>0$, $Y=1-p/\lambda>0$). Let $\pi(U,p,y)=(U,p)$.
A \emph{log-bounded entire orbit} is a solution of \eqref{eq:sys} defined and bounded on all of $\R$.
A \emph{cycle} is the image of a nonconstant periodic solution. Along a solution the Jacobian is
\[
J=\begin{pmatrix}0&1&0\\0&-\mu&-1\\x&0&-1\end{pmatrix},
\]
and the \emph{Riccati system} driven by a positive function $x(t)$ is
\begin{equation}\label{eq:ric}
\eta'=\zeta+(\mu-1)\eta-\eta^{2},\qquad \zeta'=x-(1+\eta)\zeta .
\end{equation}
For a pair $(\eta,\zeta)$ write $n=(\zeta,-\eta,-1)^{\mathsf T}$. A direct computation shows that \eqref{eq:ric} is
equivalent to
\begin{equation}\label{eq:normal}
n'+J^{\mathsf T}n+(1+\eta)n=0 .
\end{equation}
(Indeed $J^{\mathsf T}n=(-x,\ \zeta+\mu\eta,\ \eta+1)^{\mathsf T}$, so the three components of the left side of
\eqref{eq:normal} are $\zeta'-x+(1+\eta)\zeta$, $-\eta'+\zeta+(\mu-1)\eta-\eta^2$ and $0$.)

\subsection{Imported inputs}\label{sec:inputs}
\begin{enumerate}[label=(I\arabic*)]
\item \emph{Birkhoff--Hopf theorem} (Birkhoff 1957; Hopf 1963; see Seneta, \emph{Non-negative Matrices and Markov
Chains}, \S3.1, or Bushell 1973): an entrywise positive matrix $P$ contracts Hilbert's projective metric $d$ on
$\inter\R^n_+$ by the factor $\tanh(\Delta(P)/4)$, where $\Delta(P)=\sup\{d(Pu,Pw):u,w\in\R^n_+\setminus\{0\}\}$.
\item \emph{Classical planar theory} for $q'=F(q)$ with $F$ locally Lipschitz on $\R^2$ (Hartman, \emph{Ordinary
Differential Equations}, Ch.~VII, \S4, Thms 4.1--4.2 and the transversal lemmas preceding them):
\begin{enumerate}[label=(\alph*)]
\item (Poincar\'e--Bendixson) if a forward orbit is bounded and its $\omega$-limit set $\omega(q)$ contains finitely
many equilibria, then $\omega(q)$ is an equilibrium, or a periodic orbit, or consists of equilibria together with
orbits each of which converges to an equilibrium as $t\to+\infty$ and as $t\to-\infty$;
\item (monotonicity on transversals) successive intersections of an orbit with a transversal segment are monotone
along the segment, and a periodic orbit meets a transversal segment in at most one point.
\end{enumerate}
\item \emph{Invariant manifolds}: the stable/unstable manifold theorem at a hyperbolic equilibrium and at a
hyperbolic fixed point of a $C^1$ map (applied to Poincar\'e maps of hyperbolic cycles), and Lyapunov stability of
a cycle all of whose transverse Floquet multipliers have modulus $<1$ (Hartman, Ch.~IX).
\item \emph{Topology of the plane}: the Jordan curve theorem, the Schoenflies theorem, and the fact that a
homeomorphism of the circle onto a Jordan curve has winding number $\pm1$ about each interior point.
\item \emph{Whitney's extension theorem} in class $C^1$ (Whitney 1934) --- used only for the reformulation in
Theorem~\ref{thm:A}(A8)(a) at the outermost cycle.
\item From the v2.0 manuscript (\emph{Finite-parameter obstructions and cycle geometry in the Wilhelm--Heinrich
oscillator}): only its Proposition 2.1 and Theorem 4.1, and both are re-derived here in the form needed
(Lemma~\ref{lem:cyclegeom} and the first paragraph of \S\ref{sec:mult}), so the note does not depend on v2.0.
\end{enumerate}
\emph{Not needed, but consistent:} Smith (2012, Thms 1.3--1.4) proves that \eqref{eq:sys} has a compact attractor
$A_1$ in the open positive octant attracting all compact subsets of the octant, and that for $k>k_H$ every positive
orbit off the stable manifold of $E$ converges to a cycle, there being finitely many cycles (Thm 1.1(c), via the
competitive Poincar\'e--Bendixson theorem and Zhu--Smith 1994). Lemma~\ref{lem:diss} re-proves the permanence
statement directly so that the constants are explicit; the set $\cA$ below is Smith's $A_1$. The competitive
Poincar\'e--Bendixson theorem and the finiteness of cycles are \emph{not} used: on $\cA$ the classical planar theory
(I2) suffices.

\subsection{Main result}
\begin{theorem}[Attractor graph]\label{thm:A}
Let $\mu>0$, $k>0$.
\begin{enumerate}[label=\textup{(A\arabic*)}]
\item \textup{(Global attractor.)} $\cD$ is forward invariant and the semiflow on $\cD$ has a compact, connected,
invariant global attractor $\cA\subset\cD$ which attracts every compact subset of $\cD$. $\cA$ is exactly the union
of all log-bounded entire orbits of \eqref{eq:sys} in $\R^3$.
\item \textup{(Positive adjoint ratios.)} Along every log-bounded entire orbit $z(\cdot)$, with $x=ke^{U}$, the
Riccati system \eqref{eq:ric} has exactly one solution on $\R$ with $\eta\ge0$, $\zeta\ge0$. It is bounded,
\begin{equation}\label{eq:bounds}
\eta_-\le\eta\le\eta_+,\qquad \zeta_-\le\zeta\le\zeta_+,\qquad \eta_->\max(0,\mu-1),\quad \zeta_->0,
\end{equation}
with the explicit constants \eqref{eq:constants}, which depend only on $\mu$ and on
$x_-=\min_{\cA}ke^{U}$, $x_+=\max_{\cA}ke^{U}$. It attracts all nonnegative data in pull-back and forward sense at
a uniform exponential rate \eqref{eq:pullback}. The resulting functions $\eta,\zeta:\cA\to(0,\infty)$,
$a\mapsto(\eta,\zeta)(a)$ (value at time $0$ for the orbit through $a$) are continuous, and
$t\mapsto(\eta,\zeta)(\varphi_t a)$ is the Riccati solution along the orbit of $a$. On a cycle they are the
periodic Perron--Floquet ratios of v2.0; at $E$ they are the constants $\eta_E=\gamma_E-1$,
$\zeta_E=(\gamma_E-1)(\gamma_E-\mu)$, $\gamma_E(\gamma_E-1)(\gamma_E-\mu)=k$, $\gamma_E>\max(1,\mu)$.
\item \textup{(Supporting planes.)} For all $a\neq b$ in $\cA$,
\begin{equation}\label{eq:support}
y_b-y_a>\zeta(a)\,(U_b-U_a)-\eta(a)\,(p_b-p_a).
\end{equation}
\item \textup{(Convex $C^1$ graph.)} $\pi|_{\cA}$ is a homeomorphism onto the compact connected set $D=\pi(\cA)$.
The function
\[
H(u,v)=\max_{a\in\cA}\bigl\{y_a+\zeta(a)(u-U_a)-\eta(a)(v-p_a)\bigr\}
\]
is convex and globally Lipschitz on $\R^2$, and $\cA=\{(u,v,h(u,v)):(u,v)\in D\}$ with $h=H|_D$. With
$g:=(\zeta,-\eta)\circ(\pi|_\cA)^{-1}:D\to\R^2$ (continuous) one has, for all $q,q'\in D$,
\begin{equation}\label{eq:whitney}
0\le h(q')-h(q)-g(q)\cdot(q'-q)\le\bigl(g(q')-g(q)\bigr)\cdot(q'-q).
\end{equation}
Consequently $h$ is differentiable at every interior point of $D$, $h\in C^1(\inter D)$ and
\[
\nabla h=(\zeta,-\eta),\qquad h_u=\zeta>0,\qquad h_v=-\eta<-\max(0,\mu-1)\quad\text{on }\inter D,
\]
and $(h,g)$ is a Whitney $C^1$ jet on the compact set $D$.
\item \textup{(Reduced planar flow.)} The field $F(u,v)=(v,\ k-\mu v-H(u,v))$ is globally Lipschitz on $\R^2$; its
flow $\psi_t$ leaves $D$ invariant and $\pi\circ\varphi_t=\psi_t\circ\pi$ on $\cA$. Every $(u,v)\in D$ satisfies
\begin{equation}\label{eq:pde}
\zeta\,v-\eta\,(k-\mu v-h)=ke^{u}-h ,
\end{equation}
i.e.\ $h_uv+h_v(k-\mu v-h)=ke^{u}-h$ on $\inter D$. On $\inter D$ the field $F$ is $C^1$ with
$\diver F=\eta-\mu$. Moreover $K:=\zeta U-\eta p-y+k>0$ and $W:=\zeta'p-\eta'p'>0$ on $\cA\setminus\{E\}$, where
$'$ is the derivative along the flow, and $K'+(1+\eta)K=k(Ue^{U}-e^{U}+1)$, $W'+(1+\eta)W=xp^2$.
\end{enumerate}
Assume in addition $k>k_H$. Then:
\begin{enumerate}[label=\textup{(A\arabic*)},resume]
\item \textup{($E$ is a source in $\cA$.)} With $C_+=\{U\ge0,\ p\le0,\ y-k\le0\}$ and $C_-=-C_+$,
$\cA\cap(E+C_\pm)=\{E\}$. No log-bounded entire orbit other than $E$ converges to $E$ as $t\to+\infty$. The
two-dimensional unstable manifold $W^u(E)$ lies in $\cA$ and $(0,0)\in\inter D$. Every orbit in $\cA\setminus\{E\}$
has a cycle as $\omega$-limit set; its $\alpha$-limit set is $\{E\}$ or a cycle; and no non-periodic orbit has the
same cycle as $\alpha$- and $\omega$-limit set. In particular cycles exist.
\item \textup{(The domain.)} The set of crests $\{u>0:(u,0)\text{ lies on a projected cycle}\}$ is a nonempty compact
subset of $(0,\infty)$. Let $\Gamma_{\rm out}$ and $\Gamma_{\rm in}$ be the cycles with the largest and smallest
crest. Then $D$ is exactly the closed Jordan domain bounded by $\pi(\Gamma_{\rm out})$:
$\partial D=\pi(\Gamma_{\rm out})$, $\inter D$ is the bounded component of $\R^2\setminus\pi(\Gamma_{\rm out})$, and
$\cA$ is a closed topological disc with boundary $\Gamma_{\rm out}$.
\item \textup{(Multipliers.)} Let $\Gamma$ be a cycle with period $T$, $m_t=\exp\int_0^T(\eta-\mu)\,dt$,
$\lambda_c=\avg\eta-\mu$.
\begin{enumerate}[label=\textup{(\alph*)}]
\item If $\pi(\Gamma)\subset\inter D$ (every cycle except $\Gamma_{\rm out}$), then $\pi(\Gamma)$ is a periodic orbit
of the $C^1$ planar field $F$ and its planar Floquet multiplier $\exp\oint\diver F\,dt$ equals $m_t$. For
$\Gamma_{\rm out}$ the same holds for the $C^1$ field $\widehat F=(v,k-\mu v-\widehat h)$, where $\widehat h$ is any
Whitney extension of the jet $(h,g)$; $\widehat F=F$ on $D$.
\item If $m_t>1$, no orbit in $\cA\setminus\Gamma$ has $\omega$-limit set $\Gamma$. If $m_t<1$, no orbit in
$\cA\setminus\Gamma$ has $\alpha$-limit set $\Gamma$.
\item $\lambda_c(\Gamma_{\rm in})\le0$ and $\lambda_c(\Gamma_{\rm out})\le0$.
\item If every cycle has $\lambda_c<0$, there is exactly one cycle. Equivalently, if there are at least two cycles,
some cycle has $\lambda_c\ge0$.
\item If every cycle is hyperbolic ($\lambda_c\neq0$), the number of cycles is finite and odd, and ordered by
nesting they alternate $\lambda_c<0$, $\lambda_c>0$, \dots, $\lambda_c<0$.
\item \textup{(Dulac form.)} Let $\Lambda$ be a real function on the union of all cycles such that
$t\mapsto\Lambda(\varphi_ta)$ is absolutely continuous with derivative $\mathbb D\Lambda$. If
$\mathbb D\Lambda+\eta-\mu\le0$ on every cycle and is not a.e.\ zero on any cycle, the cycle is unique and
hyperbolic attracting. The reverse inequality (with the same nondegeneracy) is impossible.
\end{enumerate}
\end{enumerate}
\end{theorem}
For $0<k\le k_H$ the equilibrium is globally asymptotically stable in $\cD$ (Smith 2012, Thm 1.1(b); Boros--Hofbauer
2023, Thm 13 at $k=k_H$), so $\cA=\{E\}$ and (A1)--(A5) are trivial there.

\section{Dissipativity and the global attractor: proof of (A1)}\label{sec:diss}
Use $x=ke^{U}>0$, $q:=\lambda-p$ and $y$. Since $\mu\lambda=k$, \eqref{eq:sys} reads
\begin{equation}\label{eq:xyq}
x'=(\lambda-q)x,\qquad y'=x-y,\qquad q'=y-\mu q,
\end{equation}
and $\cD=\{x>0,\ y>0,\ q>0\}$.

\begin{lemma}[Explicit absorbing box]\label{lem:diss}
Every solution starting in $\cD$ exists for all $t\ge0$ and stays in $\cD$. There are constants
$0<\delta_x<M_x$, $0<\delta_y<M_y$, $0<\delta_q<M_q$, depending only on $(\mu,k)$, such that every solution in $\cD$
satisfies, for all sufficiently large $t$,
\[
(x,y,q)(t)\in B_0:=[\delta_x,M_x]\times[\delta_y,M_y]\times[\delta_q,M_q].
\]
\end{lemma}
\begin{proof}
On the maximal interval of existence $x(t)=x(0)\exp\int_0^t(\lambda-q)>0$ and, by variation of constants,
\begin{equation}\label{eq:voc}
y(t)=e^{-(t-s)}y(s)+\int_s^te^{-(t-r)}x(r)\,dr,\qquad q(t)=e^{-\mu(t-s)}q(s)+\int_s^te^{-\mu(t-r)}y(r)\,dr ,
\end{equation}
so $y,q>0$ persist. As long as $q>0$ we have $x'\le\lambda x$, so $x\le x(0)e^{\lambda t}$ and then $y,q$ are bounded
on bounded intervals by \eqref{eq:voc}. Hence solutions are global and $\cD$ is forward invariant.

\emph{Upper bounds.} From $x'/x\le\lambda$, $x(r)\ge x(t)e^{-\lambda(t-r)}$ for $0\le r\le t$. Dropping the
nonnegative initial terms in \eqref{eq:voc} (with $s=0$) twice,
\[
q(t)\ge\int_0^te^{-\mu(t-s)}\!\!\int_0^se^{-(s-r)}x(r)\,dr\,ds\ \ge\ c(t)\,x(t),\qquad
c(t):=\iint_{a,b\ge0,\ a+b\le t}e^{-\mu a}e^{-b}e^{-\lambda(a+b)}\,da\,db .
\]
The function $c$ is positive and nondecreasing; put $c_1=c(1)$. For $t\ge1$, $x'\le x(\lambda-c_1x)$, and comparison
with the logistic equation gives $\limsup x\le\lambda/c_1$. Put $M_x=2\lambda/c_1$; then $x\le M_x$ for $t\ge t_1$.
By \eqref{eq:voc}, $y(t)\le e^{-(t-t_1)}y(t_1)+M_x$, so $y\le M_y:=2M_x$ for $t\ge t_2$, and likewise
$q\le M_q:=2M_y/\mu$ for $t\ge t_3$.

\emph{Lower bounds.} Let $m=\min(1,\mu)$, $a=k/4$ (so $a/\mu=\lambda/4$), and choose $L>0$ with
$(M_q+M_ys)e^{-ms}\le\lambda/4$ for all $s\ge L$. Put $\delta_x=a\,e^{-M_qL}$. Consider $t_0\ge t_3$ and an interval
$[t_0,t_0+\sigma]$, $0<\sigma\le\infty$, on which $x\le a$. For $0\le s\le\sigma$, \eqref{eq:voc} gives
$y(t_0+s)\le M_ye^{-s}+a$ and then
\[
q(t_0+s)\le M_qe^{-\mu s}+M_y\int_0^se^{-\mu(s-r)}e^{-r}dr+\frac a\mu\le(M_q+M_ys)e^{-ms}+\frac\lambda4 ,
\]
because $\int_0^se^{-\mu(s-r)}e^{-r}dr\le\int_0^se^{-m(s-r)}e^{-mr}dr=se^{-ms}$. Hence $q\le\lambda/2$ and
$x'\ge(\lambda/2)x$ on $[t_0+L,t_0+\sigma]$. Consequently $\sigma<\infty$ (otherwise $x\to\infty$, contradicting
$x\le a$): every maximal interval in $[t_3,\infty)$ on which $x\le a$ is bounded. Let $t_4\ge t_3$ be a time with
$x(t_4)\ge a$, and let $[t_0,t_0+\sigma]\subset[t_4,\infty)$ be a maximal interval with $x\le a$; then $x(t_0)=a$.
On $[t_0,t_0+\min(\sigma,L)]$, $x'/x=\lambda-q\ge-M_q$ gives $x\ge ae^{-M_qL}=\delta_x$; on the rest of the interval
$x$ is nondecreasing. So $x\ge\delta_x$ on $[t_4,\infty)$. Then \eqref{eq:voc} gives
$y(t)\ge\delta_x(1-e^{-(t-t_4)})\ge\delta_y:=\delta_x/2$ eventually and then $q\ge\delta_q:=\delta_y/(2\mu)$
eventually.
\end{proof}
In the coordinates $(U,p,y)$, $B_0$ is the compact box
$[\log(\delta_x/k),\log(M_x/k)]\times[\lambda-M_q,\lambda-\delta_q]\times[\delta_y,M_y]\subset\cD$.

\begin{proof}[Proof of (A1)]
Let $N\subset\cD$ be a compact box (in the coordinates $(U,p,y)$) whose interior contains $B_0$.

\emph{(a) A compact forward-invariant absorbing set.} For $z\in\cD$ Lemma~\ref{lem:diss} gives $t_z\ge1$ with
$\varphi_{t_z}(z)\in B_0\subset\inter N$, and by continuity an open $V_z\ni z$ with
$\varphi_{t_z}(V_z)\subset\inter N$. Cover $N$ by finitely many $V_{z_i}$, let $\tau=\max t_{z_i}$, and put
$B_1=\varphi([0,\tau]\times N)$, a compact subset of $\cD$. If $z\in N$ then $z\in V_{z_i}$ for some $i$,
$\varphi_t(z)\in B_1$ for $0\le t\le t_{z_i}$ and $\varphi_{t_{z_i}}(z)\in N$; since $t_{z_i}\ge1$, induction gives
$\varphi_t(z)\in B_1$ for all $t\ge0$. Hence $B_1$ is forward invariant. Any compact $C\subset\cD$ is covered by
finitely many $V_z$, so $\varphi_t(C)\subset B_1$ for all $t\ge t_C$.

\emph{(b) The attractor.} Put $\cA=\bigcap_{t\ge0}\varphi_t(B_1)$, a decreasing intersection of nonempty compact
sets. Then $\operatorname{dist}(\varphi_t(B_1),\cA)\to0$, so $\cA$ attracts $B_1$ and, by (a), every compact subset
of $\cD$. Invariance: if $z\in\cA$ and $s\ge0$ then $\varphi_s(z)\in\varphi_{s+t}(B_1)$ for all $t\ge0$ and
$\varphi_s(z)\in\varphi_s(B_1)\subset\varphi_{t'}(B_1)$ for $t'\le s$, so $\varphi_s(\cA)\subset\cA$. Conversely,
for $z\in\cA$ and $n\ge s$ write $z=\varphi_n(b_n)=\varphi_s(\varphi_{n-s}(b_n))$ with $b_n\in B_1$; by injectivity
of $\varphi_s$ the point $w=\varphi_{n-s}(b_n)$ does not depend on $n$, lies in $\varphi_{n-s}(B_1)$ for all
$n\ge s$, hence in $\cA$, and $\varphi_s(w)=z$. So $\varphi_s(\cA)=\cA$, and each point of $\cA$ lies on an entire
orbit contained in the compact set $\cA$.

\emph{(c) Characterisation.} Let $z(\cdot)$ be a solution of \eqref{eq:sys} on $\R$ that is bounded. Then
$0<x_-\le x(t)\le x_+$. The bounded function $y$ solves $y'=x-y$ on $\R$, and two bounded solutions of this equation
differ by $ce^{-t}$ with $c=0$; hence $y(t)=\int_{-\infty}^te^{-(t-s)}x(s)\,ds\ge x_-$. Likewise
$q=\lambda-p=\int_{-\infty}^te^{-\mu(t-s)}y(s)\,ds\ge x_-/\mu$. So the closure $C$ of the orbit is a compact subset of
$\cD$, and for every $s$ and every $t\ge0$, $z(s)=\varphi_t\bigl(\varphi_{t_C}(z(s-t-t_C))\bigr)\in\varphi_t(B_1)$.
Thus $z(s)\in\cA$. With (b) this proves the characterisation.

\emph{(d) Connectedness.} Let $N'\subset\cD$ be a compact box containing $B_1$. Then $\varphi_t(N')$ is connected,
contains $\cA=\varphi_t(\cA)$, and lies in any prescribed neighbourhood of $\cA$ for large $t$. If $\cA$ were the
union of two disjoint nonempty compact sets, separated by disjoint open sets $O_1,O_2$, then for large $t$ the
connected set $\varphi_t(N')\subset O_1\cup O_2$ would meet both. Contradiction.
\end{proof}

\section{The positive adjoint direction: proof of (A2)}\label{sec:adjoint}
Throughout this section $x:\R\to[x_-,x_+]$ is a continuous function with $0<x_-\le x_+$ (no dynamics is assumed).
Consider
\begin{equation}\label{eq:lin}
\mathsf H'=B(t)\mathsf H,\qquad B(t)=\begin{pmatrix}0&0&x(t)\\1&\mu&0\\0&1&1\end{pmatrix},
\end{equation}
with evolution operator $\Phi(t,s)$. If $\mathsf H$ solves \eqref{eq:lin} with $\mathsf H_3>0$, then
$\eta=\mathsf H_2/\mathsf H_3$, $\zeta=\mathsf H_1/\mathsf H_3$ solve \eqref{eq:ric}; conversely, if $(\eta,\zeta)$
solves \eqref{eq:ric} on an interval, then $\mathsf H_3=\exp\int(1+\eta)$, $\mathsf H_2=\eta\mathsf H_3$,
$\mathsf H_1=\zeta\mathsf H_3$ solves \eqref{eq:lin} (direct differentiation). Let $B_\pm$ be $B$ with $x$ replaced
by $x_\pm$, and
\[
m_0=\min_{i,j}\bigl(e^{B_-}\bigr)_{ij},\quad M_0=\max_{i,j}\bigl(e^{B_+}\bigr)_{ij},\quad
\Delta_0=2\log\frac{M_0}{m_0},\quad \tau_0=\tanh\frac{\Delta_0}{4}<1 .
\]
Let $d(u,w)=\log\bigl(\max_i\frac{u_i}{w_i}\,\max_i\frac{w_i}{u_i}\bigr)$ be Hilbert's projective metric on
$\inter\R^3_+$.

\begin{lemma}[Uniform positivity and contraction]\label{lem:cone}
\begin{enumerate}[label=\textup{(\roman*)}]
\item $m_0>0$, and for every $t$ all entries of $\Phi(t+1,t)$ lie in $[m_0,M_0]$. For $t\ge s$, $\Phi(t,s)$ is
entrywise nonnegative.
\item For $t-s\ge1$ and $u,w\in\R^3_+\setminus\{0\}$:
$d\bigl(\Phi(t,s)u,\Phi(t,s)w\bigr)\le\Delta_0\,\tau_0^{\lfloor t-s\rfloor-1}$.
\end{enumerate}
\end{lemma}
\begin{proof}
(i) All entries of $B(t)$, $B_\pm$ are nonnegative (the diagonal is $0,\mu,1$) and $B_-\le B(t)\le B_+$ entrywise.
The Peano--Baker series $\Phi(t,s)=I+\int_s^tB+\int_s^tB(t_1)\int_s^{t_1}B(t_2)+\cdots$ has nonnegative terms and is
entrywise monotone in $B$, so $e^{(t-s)B_-}\le\Phi(t,s)\le e^{(t-s)B_+}$. The digraph of $B_-$ contains the cycle
$1\to2\to3\to1$ (entries $(2,1)$, $(3,2)$, $(1,3)$ are positive), so $B_-$ is irreducible and
$e^{B_-}=\sum_nB_-^n/n!$ is entrywise positive.

(ii) Let $P$ have entries in $[m_0,M_0]$ and $u,w\in\R^3_+\setminus\{0\}$. Then
$(Pu)_i/(Pw)_i\in\bigl[\tfrac{m_0|u|_1}{M_0|w|_1},\tfrac{M_0|u|_1}{m_0|w|_1}\bigr]$ for each $i$, hence
$d(Pu,Pw)\le\Delta_0$ and $\Delta(P)\le\Delta_0$; by (I1), $P$ contracts $d$ by the factor $\tau_0$. Put
$n=\lfloor t-s\rfloor\ge1$ and factor
$\Phi(t,s)=\Phi(t,t-1)\cdots\Phi(t-n+2,t-n+1)\,\Phi(t-n+1,t-n)\,\Phi(t-n,s)$. The last factor maps
$\R^3_+\setminus\{0\}$ into itself (nonnegative and invertible), the next one maps it into $\inter\R^3_+$ with image
of $d$-diameter $\le\Delta_0$, and each of the remaining $n-1$ factors contracts $d$ by $\tau_0$.
\end{proof}

\begin{lemma}[Suprema of bounded $C^1$ functions]\label{lem:sup}
If $\phi\in C^1(\R)$ is bounded, there are $t_n$ with $\phi(t_n)\to\sup\phi$ and $\phi'(t_n)\to0$ (and similarly
for $\inf\phi$).
\end{lemma}
\begin{proof}
Let $S=\sup\phi$, pick $\tau_n$ with $\phi(\tau_n)\ge S-n^{-2}$ and let $t_n$ be a maximum point of
$\phi(t)-(t-\tau_n)^2/n$. Then $\phi(t_n)-(t_n-\tau_n)^2/n\ge\phi(\tau_n)\ge S-n^{-2}$, whence
$\phi(t_n)\ge S-n^{-2}$, $(t_n-\tau_n)^2\le n^{-1}$ and $|\phi'(t_n)|=2|t_n-\tau_n|/n\le2n^{-3/2}$.
\end{proof}

\begin{proposition}[Unique positive entire Riccati solution]\label{prop:riccati}
\begin{enumerate}[label=\textup{(\roman*)}]
\item There is exactly one solution $(\eta,\zeta)$ of \eqref{eq:ric} on $\R$ with $\eta\ge0$ and $\zeta\ge0$
everywhere.
\item It satisfies \eqref{eq:bounds} with
\begin{equation}\label{eq:constants}
\eta_+=\tfrac{\mu-1+\sqrt{(\mu-1)^2+4x_+}}2,\quad \zeta_-=\tfrac{x_-}{1+\eta_+},\quad
\eta_-=\tfrac{\mu-1+\sqrt{(\mu-1)^2+4\zeta_-}}2,\quad \zeta_+=\tfrac{x_+}{1+\eta_-}.
\end{equation}
\item For every $s\in\R$ and $(\eta_0,\zeta_0)\in[0,\infty)^2$ the solution $(\tilde\eta,\tilde\zeta)$ of
\eqref{eq:ric} with data $(\eta_0,\zeta_0)$ at time $s$ exists on $[s,\infty)$, and for $t\ge s+1$
\begin{equation}\label{eq:pullback}
\bigl|\log\tilde\eta(t)-\log\eta(t)\bigr|\le\Delta_0\tau_0^{\lfloor t-s\rfloor-1},\qquad
\bigl|\log\tilde\zeta(t)-\log\zeta(t)\bigr|\le\Delta_0\tau_0^{\lfloor t-s\rfloor-1}.
\end{equation}
\end{enumerate}
\end{proposition}
\begin{proof}
\emph{Existence.} Let $\Sigma=\{u\in\R^3_+:u_1+u_2+u_3=1\}$ and, for $s<t$,
$C(t,s)=\{\Phi(t,s)u/|\Phi(t,s)u|_1:u\in\Sigma\}$, a compact subset of $\Sigma$. Since
$\Phi(t,s')=\Phi(t,s)\Phi(s,s')$ and $\Phi(s,s')$ maps $\R^3_+\setminus\{0\}$ into itself, $C(t,s')\subset C(t,s)$
for $s'<s$. For $t-s\ge1$, $C(t,s)\subset\inter\Sigma$ and its $d$-diameter is at most
$\Delta_0\tau_0^{\lfloor t-s\rfloor-1}$ by Lemma~\ref{lem:cone}. Hence $\bigcap_{s<t}C(t,s)$ is nonempty (nested
compacta) and of $d$-diameter $0$, i.e.\ a single point $\hat{\mathsf h}(t)\in\inter\Sigma$ ($d$ is a metric on
$\inter\Sigma$). For $t'>t$, $\Phi(t',t)\hat{\mathsf h}(t)$, normalised, lies in $C(t',s)$ for every $s<t$, hence in
$\bigcap_{s<t}C(t',s)=\bigcap_{s<t'}C(t',s)=\{\hat{\mathsf h}(t')\}$. Therefore
$\mathsf H^*(t):=\Phi(t,0)\hat{\mathsf h}(0)$ is, for every $t\in\R$, a positive multiple of $\hat{\mathsf h}(t)$:
a positive entire solution of \eqref{eq:lin}. Its ratios $(\eta,\zeta)$ solve \eqref{eq:ric} on $\R$ and are
positive.

\emph{Uniqueness.} Let $(\tilde\eta,\tilde\zeta)$ be a solution on $\R$ with $\tilde\eta,\tilde\zeta\ge0$. The
associated $\tilde{\mathsf H}$ is a solution of \eqref{eq:lin} with values in $\R^3_+\setminus\{0\}$, so
$\tilde{\mathsf H}(t)=\Phi(t,s)\tilde{\mathsf H}(s)$ has normalisation in $C(t,s)$ for all $s<t$, i.e.\ equal to
$\hat{\mathsf h}(t)$. Hence $(\tilde\eta,\tilde\zeta)=(\eta,\zeta)$.

\emph{Bounds.} $\mathsf H^*(t)=\Phi(t,t-1)c$ with $c\in\R^3_+\setminus\{0\}$, so
$m_0|c|_1\le\mathsf H^*_i(t)\le M_0|c|_1$ and $\eta,\zeta\in[m_0/M_0,M_0/m_0]$: both are bounded and bounded away
from $0$; by \eqref{eq:ric} so are their derivatives. Write $\bar\phi=\sup\phi$, $\underline\phi=\inf\phi$.
Lemma~\ref{lem:sup} applied to $\zeta$ gives $t_n$ with $x(t_n)-(1+\eta(t_n))\zeta(t_n)\to0$, $\zeta(t_n)\to\bar\zeta$,
so $\bar\zeta\le x_+/(1+\underline\eta)\le x_+$. Applied to $\eta$ at its supremum:
$\bar\eta(\bar\eta+1-\mu)=\lim\zeta(t_n)\le x_+$ and $\bar\eta>0$, so $\bar\eta\le\eta_+$. At the infimum of
$\zeta$: $\underline\zeta=\lim x(t_n)/(1+\eta(t_n))\ge x_-/(1+\eta_+)=\zeta_-$. At the infimum of $\eta$:
$\underline\eta(\underline\eta+1-\mu)=\lim\zeta(t_n)\ge\zeta_->0$ with $\underline\eta>0$; the quadratic
$s(s+1-\mu)$ is $\le0$ on $[0,\max(0,\mu-1)]$ and increasing beyond, so $\underline\eta\ge\eta_->\max(0,\mu-1)$.
Finally $\bar\zeta\le x_+/(1+\eta_-)=\zeta_+$.

\emph{Attraction.} The solution of \eqref{eq:lin} with $\tilde{\mathsf H}(s)=(\zeta_0,\eta_0,1)^{\mathsf T}$ is
nonnegative, and $\tilde{\mathsf H}_3'=\tilde{\mathsf H}_2+\tilde{\mathsf H}_3\ge\tilde{\mathsf H}_3$ keeps
$\tilde{\mathsf H}_3>0$; its ratios are the (unique, since \eqref{eq:ric} is polynomial) solution of \eqref{eq:ric}
with the given data, defined on $[s,\infty)$. By Lemma~\ref{lem:cone}(ii) with $u=\tilde{\mathsf H}(s)$,
$w=\mathsf H^*(s)$, $d(\tilde{\mathsf H}(t),\mathsf H^*(t))\le\Delta_0\tau_0^{\lfloor t-s\rfloor-1}$, and
$|\log(u_i/u_3)-\log(w_i/w_3)|=|\log(u_i/w_i)-\log(u_3/w_3)|\le d(u,w)$.
\end{proof}

\begin{proof}[Proof of (A2)]
For $a\in\cA$ put $x_a(t)=k\exp U(\varphi_ta)$, $t\in\R$; then $x_-\le x_a\le x_+$ with $x_\pm$ the extreme values of
$ke^{U}$ on the compact set $\cA$. Proposition~\ref{prop:riccati} gives the unique nonnegative entire solution
$(\eta_a,\zeta_a)(t)$, the bounds, and \eqref{eq:pullback} with constants uniform in $a$. Define
$(\eta,\zeta)(a)=(\eta_a,\zeta_a)(0)$. Since $x_{\varphi_sa}(t)=x_a(t+s)$, uniqueness gives
$(\eta,\zeta)(\varphi_sa)=(\eta_a,\zeta_a)(s)$.

\emph{Continuity.} For $s\ge1$ let $R_s(a)$ be the value at time $0$ of the solution of \eqref{eq:ric}, driven by
$x_a$, with data $(1,1)$ at time $-s$. The flow restricted to the compact invariant set $\cA$ is a continuous action
of $\R$, so $a\mapsto x_a|_{[-s,0]}$ is continuous into $C([-s,0])$ and, by continuous dependence, $R_s:\cA\to
(0,\infty)^2$ is continuous. By \eqref{eq:pullback}, $\sup_{a\in\cA}|\log R_s(a)-\log(\eta,\zeta)(a)|\le
\Delta_0\tau_0^{\lfloor s\rfloor-1}\to0$. A uniform limit of continuous functions is continuous.

\emph{Cycles and $E$.} On a cycle, $(\eta,\zeta)(\cdot+T)$ is again a nonnegative entire solution, so
$(\eta,\zeta)$ is $T$-periodic; the Perron--Floquet ratios of v2.0 are positive entire solutions and hence coincide
with it. At $E$ the constant solutions of \eqref{eq:ric} with $x=k$ satisfy $\zeta=k/(1+\eta)$,
$\eta(\eta+1-\mu)=\zeta$, i.e.\ $\gamma(\gamma-1)(\gamma-\mu)=k$ with $\gamma=1+\eta$; positivity of $\eta,\zeta$
forces $\gamma>\max(1,\mu)$, where the cubic increases from $0$ to $\infty$.
\end{proof}

\section{Supporting planes: proof of (A3)}\label{sec:support}
\begin{theorem}[Two-orbit comparison]\label{thm:support}
Let $z_1,z_2$ be log-bounded entire solutions of \eqref{eq:sys} with $z_1\not\equiv z_2$, and let $(\eta,\zeta)$ be
the positive entire Riccati solution along $z_1$, $n=(\zeta,-\eta,-1)^{\mathsf T}$. Then
$\Delta:=n\cdot(z_2-z_1)=\zeta(U_2-U_1)-\eta(p_2-p_1)-(y_2-y_1)$ satisfies
\begin{equation}\label{eq:deltaode}
\Delta'+(1+\eta)\Delta=-x_1\bigl[e^{U_2-U_1}-1-(U_2-U_1)\bigr]=:-r(t)\le0
\end{equation}
and $\Delta(t)<0$ for all $t\in\R$.
\end{theorem}
\begin{proof}
With $d=z_2-z_1$ one has $d'=f(z_2)-f(z_1)=J_1d+(0,0,r)^{\mathsf T}$, because the only nonlinear term of $f$ is
$ke^{U}$ in the third component and $ke^{U_2}-ke^{U_1}=x_1(U_2-U_1)+r$. By \eqref{eq:normal},
$\Delta'=n'\cdot d+n\cdot J_1d+n_3r=-(1+\eta)\Delta-r$. For $s<t$,
\[
\Delta(t)=\Delta(s)e^{-\int_s^t(1+\eta)}-\int_s^te^{-\int_\tau^t(1+\eta)}r(\tau)\,d\tau .
\]
$\Delta$ is bounded and $1+\eta\ge1$; letting $s\to-\infty$ gives
$\Delta(t)=-\int_{-\infty}^te^{-\int_\tau^t(1+\eta)}r(\tau)\,d\tau\le0$. If $\Delta(t_0)=0$, then $r\equiv0$ on
$(-\infty,t_0]$; as $x_1>0$ and $e^{s}-1-s>0$ for $s\ne0$, $U_2\equiv U_1$ there, hence $p_2=U_2'=p_1$ and
$y_2=k-\mu p_2-p_2'=y_1$ there, and $z_2\equiv z_1$ by uniqueness. Contradiction.
\end{proof}
\begin{proof}[Proof of (A3)]
For $a\ne b$ in $\cA$ apply Theorem~\ref{thm:support} to $z_1=\varphi_{\cdot}(a)$, $z_2=\varphi_{\cdot}(b)$ (distinct
as functions, also when $b$ lies on the orbit of $a$) at $t=0$.
\end{proof}

\begin{lemma}[Gauges]\label{lem:gauges}
On $\cA$: \textup{(i)} $O:=n\cdot f=\zeta p-\eta(k-\mu p-y)-(ke^{U}-y)=0$;
\textup{(ii)} $K=\zeta U-\eta p-y+k>0$ off $E$, and $K'+(1+\eta)K=k(Ue^{U}-e^{U}+1)$;
\textup{(iii)} $W=\zeta'p-\eta'p'>0$ off $E$, and $W'+(1+\eta)W=xp^2$;
\textup{(iv)} $\zeta_EU-\eta_Ep-(y-k)<0$ off $E$.
\end{lemma}
\begin{proof}
(i) Along an orbit in $\cA$, $z''=Jz'$, so by \eqref{eq:normal} $O'=(n'+J^{\mathsf T}n)\cdot z'=-(1+\eta)O$. $O$ is
bounded on $\R$ and $O(t)=O(s)e^{-\int_s^t(1+\eta)}\to0$ as $s\to-\infty$. (ii) $K(a)=-n(a)\cdot(E-a)>0$ by
\eqref{eq:support} with $b=E$, and the equation is \eqref{eq:deltaode} with $z_2\equiv E$:
$r=x(e^{-U}-1+U)=k(1-e^{U}+Ue^{U})$. (iv) is \eqref{eq:support} with $a=E$. (iii) Using \eqref{eq:ric}, $p''=-\mu
p'-(x-y)$ and (i) in the form $x-y=\zeta p-\eta p'$, i.e.\ $p''=(\eta-\mu)p'-\zeta p$:
\begin{align*}
W'&=\zeta''p+\zeta'p'-\eta''p'-\eta'p''\\
&=\bigl(xp-\eta'\zeta-(1+\eta)\zeta'\bigr)p-\bigl((\mu-1)\eta'-2\eta\eta'\bigr)p'-\eta'\bigl((\eta-\mu)p'-\zeta p\bigr)\\
&=xp^2-(1+\eta)\zeta'p+(1+\eta)\eta'p'=xp^2-(1+\eta)W .
\end{align*}
$W$ is bounded on $\R$, so $W(t)=\int_{-\infty}^te^{-\int_\tau^t(1+\eta)}xp^2\,d\tau\ge0$, with equality only if
$p\equiv0$ on $(-\infty,t]$, which forces $y\equiv k$, $x\equiv k$: the orbit is $E$.
\end{proof}

\section{The graph and the reduced flow: proof of (A4), (A5)}\label{sec:graph}
\begin{proof}[Proof of (A4)]
By (A2), $a\mapsto(y_a,\zeta(a),\eta(a),\pi(a))$ is continuous on the compact set $\cA$; so $H$ is a maximum of
affine functions with slopes in a compact set: convex, finite, Lipschitz with constant
$\max_\cA|(\zeta,\eta)|$. For $b\in\cA$ the term $a=b$ equals $y_b$ at $\pi(b)$ and every other term is $<y_b$ there
by \eqref{eq:support}; so $H(\pi(b))=y_b$. Hence $\pi|_\cA$ is injective, $\cA$ is the graph of $h=H|_D$, and
$\pi|_\cA$, a continuous bijection of a compact space onto $D$, is a homeomorphism; $D$ is compact and connected by
(A1). So $g=(\zeta,-\eta)\circ(\pi|_\cA)^{-1}$ is continuous on $D$. For $q=\pi(a)$, $q'=\pi(b)$, \eqref{eq:support}
reads $h(q')-h(q)\ge g(q)\cdot(q'-q)$; exchanging $a$ and $b$, $h(q)-h(q')\ge g(q')\cdot(q-q')$. Together these are
\eqref{eq:whitney}. Therefore, with $\omega_g$ the modulus of continuity of $g$ on $D$,
\[
\bigl|h(q')-h(q)-g(q)\cdot(q'-q)\bigr|\le\omega_g(|q'-q|)\,|q'-q|\qquad(q,q'\in D),
\]
which is Whitney's $C^1$ condition for the jet $(h,g)$ on $D$, and which at an interior point $q$ is
differentiability of $h$ with $\nabla h(q)=g(q)$. Continuity of $g$ gives $h\in C^1(\inter D)$; the sign statements
are \eqref{eq:bounds}.
\end{proof}
\begin{proof}[Proof of (A5)]
$F$ is globally Lipschitz because $H$ is, so $\psi_t$ is a complete flow on $\R^2$. For $q\in D$ let
$a=(q,h(q))\in\cA$ and $z(t)=\varphi_t(a)\in\cA$. Then $y(t)=H(U(t),p(t))$, so $(U,p)$ solves $q'=F(q)$; by
uniqueness $\pi(\varphi_ta)=\psi_t(q)$, and $D$ is invariant. Equation \eqref{eq:pde} is Lemma~\ref{lem:gauges}(i)
written at the point $(u,v,h(u,v))$; on $\inter D$, $(\zeta,-\eta)=\nabla h$. On $\inter D$, $F=(v,k-\mu v-h)$ with
$h\in C^1$, and $\diver F=\partial_v(k-\mu v-h)=-\mu+\eta$. The gauges are Lemma~\ref{lem:gauges}.
\end{proof}
\begin{remark}\label{rem:convexgauges}
In terms of $h$: $K=k-[h-uh_u-vh_v]$ is the height of $E$ above the tangent plane at $(u,v)$, and
$W=\frac{d}{dt}\bigl(\nabla h(\psi_tq)\bigr)\cdot F$ is the second derivative of $h$ along the flow line (when it
exists, $W=F^{\mathsf T}D^2h\,F$). Both are convexity gauges. Also, on $D\cap\{v=0\}$, \eqref{eq:pde} gives
$h(u,0)-k=k(e^{u}-1)/(1+\eta(u,0))$, so $v'=-k(e^{u}-1)/(1+\eta)$ there: the reduced flow crosses the positive
$u$-axis downwards and the negative $u$-axis upwards, and $(0,0)$ is the only equilibrium of $F$ in $D$ (if
$F(q)=0$, $q\in D$, then $\psi_tq=q$, so $(q,h(q))$ is an equilibrium of \eqref{eq:sys}, i.e.\ $E$).
\end{remark}

\section{Above Hopf: proof of (A6), (A7)}\label{sec:above}
Assume $k>k_H$. The characteristic polynomial of $J_E$ is $z^3+Az^2+\mu z+k$; it has the real root
$-\gamma_E<-A$ (the cubic $\gamma(\gamma-1)(\gamma-\mu)$ equals $k_H<k$ at $\gamma=A$ and increases) and, since the
remaining two roots have sum $\gamma_E-A>0$ and product $k/\gamma_E>0$ and cannot be nonnegative reals (positive
coefficients), a pair with positive real part. So $E$ is hyperbolic with a one-dimensional stable manifold tangent
to the eigenvector $v_s=(1,-\gamma_E,-\gamma_E(\gamma_E-\mu))^{\mathsf T}\in\inter C_+$ and a two-dimensional
unstable manifold tangent to the real plane $E^u$. By \eqref{eq:normal} with constant $n_E=(\zeta_E,-\eta_E,-1)$,
$J_E^{\mathsf T}n_E=-\gamma_En_E$; a left eigenvector for the simple eigenvalue $-\gamma_E$ annihilates the other
eigenvectors, so
\begin{equation}\label{eq:Eu}
E^u=n_E^{\perp}=\{\delta y=\zeta_E\,\delta U-\eta_E\,\delta p\},
\end{equation}
a graph over the $(U,p)$-plane.

\begin{proof}[Proof of (A6)]
\emph{Cone exclusion.} Let $a\in\cA\setminus\{E\}$ and $r=a-E=(U,p,y-k)$. If $r\in C_+$ then
$\zeta_EU-\eta_Ep-(y-k)\ge0$, contradicting Lemma~\ref{lem:gauges}(iv). If $r\in C_-$ then
$K=\zeta U-\eta p-(y-k)\le0$ (as $\zeta,\eta>0$), contradicting Lemma~\ref{lem:gauges}(ii).

\emph{No forward convergence to $E$.} If $a\in\cA\setminus\{E\}$ and $\varphi_t(a)\to E$, then by the stable manifold
theorem $\varphi_t(a)$ lies on the local stable manifold for large $t$, a $C^1$ curve tangent to $v_s$; hence
$(\varphi_t(a)-E)/|\varphi_t(a)-E|\to\pm v_s/|v_s|\in\inter C_+\cup\inter C_-$ and $\varphi_t(a)\in E+C_\pm$ for
large $t$, contradicting cone exclusion.

\emph{$W^u(E)\subset\cA$ and $(0,0)\in\inter D$.} A point of $W^u(E)$ has a backward orbit converging to $E$ (hence
bounded and, near $E$, inside $\cD$) and a forward orbit which stays in $\cD$ and is bounded by
Lemma~\ref{lem:diss}; by (A1) it lies in $\cA$. The local unstable manifold is a $C^1$ surface tangent at $E$ to
$E^u$, on which $d\pi$ is injective by \eqref{eq:Eu}; by the inverse function theorem $\pi(W^u_{\rm loc}(E))$
contains a neighbourhood $N_E$ of $(0,0)$ in $\R^2$. So $N_E\subset D$, and every point of $N_E\setminus\{0\}$ has
$\alpha$-limit set $\{0\}$ under $\psi$; in particular $N_E$ contains no periodic point.

\emph{Limit sets.} Let $q\in D\setminus\{0\}$. Its $\psi$-orbit lies in the compact invariant set $D$, in which $0$
is the only equilibrium of $F$ (Remark~\ref{rem:convexgauges}). By (I2a), $\omega(q)$ is $\{0\}$, or a periodic
orbit, or contains an orbit $\ne\{0\}$ in $D$ converging to $0$ as $t\to+\infty$. The first and third alternatives
contradict the previous step (lift to $\cA$ by (A5)). So $\omega(q)$ is a periodic orbit of $\psi$ in $D$, i.e.\ the
projection of a cycle. Applying (I2a) to $-F$: $\alpha(q)$ is $\{0\}$, or a periodic orbit, or contains an orbit
$\ne\{0\}$ in $D$ which converges to $0$ as $t\to\pm\infty$ --- again excluded.

\emph{$\alpha(q)\neq\omega(q)$ for non-periodic $q$.} Suppose $\alpha(q)=\omega(q)=J$ periodic, $q\notin J$. Take a
transversal segment $S$ through a point $c\in J$. The orbit of $q$ lies in one component of $\R^2\setminus J$, so
it meets $S$ only on one side of $c$; it meets $S$ at times $\to+\infty$ and at times $\to-\infty$ in points
converging to $c$. By (I2b) these points, ordered by time, form a strictly monotone bi-infinite sequence on one
side of $c$ (strict because $q$ is not periodic), which cannot converge to $c$ at both ends.

\emph{Existence.} For $q\in N_E\setminus\{0\}$, $\omega(q)$ is the projection of a cycle.
\end{proof}

\begin{lemma}[Geometry of projected cycles; cf.\ v2.0, Theorem 4.1]\label{lem:cyclegeom}
Let $J\subset D$ be a periodic orbit of $\psi$ (equivalently, by (A5), the projection of a cycle). Then $J$ is a
Jordan curve with $0$ in its interior, traversed clockwise, and it meets the positive $u$-axis in exactly one point
$(c,0)$, its \emph{crest}, and the negative $u$-axis in exactly one point. The points $(u,0)$ with $0<u<c$ lie inside
$J$ and those with $u>c$ outside. Two such orbits with crests $c<c'$ are disjoint and satisfy
$J\subset\operatorname{int}J'$, $J'\subset\operatorname{ext}J$.
\end{lemma}
\begin{proof}
$J$ is a Jordan curve not through the equilibrium $0$. Since $u'=v$, the periodic function $u$ has extrema, where
$v=0$. By Remark~\ref{rem:convexgauges}, on $D\cap\{v=0\}$ one has $v'=-k(e^{u}-1)/(1+\eta)$: every crossing of
$\{v=0,u>0\}$ is transverse and downward, every crossing of $\{v=0,u<0\}$ transverse and upward, and $J$ does not
meet $\{v=0,u=0\}$. At the maximum of $u$ along $J$, $v=0$ and $v'\le0$, so $u_{\max}>0$; likewise $u_{\min}<0$. The
winding number of $J$ about $0$ is the signed number of its crossings of the ray $\{v=0,u>0\}$; it is $0$ or $\pm1$
for a Jordan curve, and all crossings have the same sign and there is at least one; so there is exactly one, the
winding number is $-1$ (clockwise) and $0$ is inside $J$. The same holds for the negative ray. The open segment from
$0$ to $(c,0)$ does not meet $J$ and contains points near $0$, so it lies inside $J$; the ray beyond $(c,0)$ does not
meet $J$ and is unbounded, so it lies outside. Distinct periodic orbits are disjoint; if $c<c'$, the point $(c,0)\in J$
is inside $J'$ and the point $(c',0)\in J'$ is outside $J$, so by connectedness $J\subset\operatorname{int}J'$ and
$J'\subset\operatorname{ext}J$.
\end{proof}
Write $J_c$ for the projected cycle with crest $c$, and let $\mathrm{Cr}\subset(0,\infty)$ be the set of crests.

\begin{lemma}\label{lem:crests}
$\mathrm{Cr}$ is nonempty and compact in $(0,\infty)$.
\end{lemma}
\begin{proof}
Nonempty by (A6); bounded since $D$ is; bounded away from $0$ since $N_E$ contains no periodic point. Let
$u_n\in\mathrm{Cr}$, $u_n\to\bar u$, and $\bar q=(\bar u,0)\in D\setminus\{0\}$. Suppose $\bar q$ is not periodic. By
(A6), $\omega(\bar q)=J_{c_\omega}$ and $\alpha(\bar q)$ is $\{0\}$ (put $c_\alpha=0$, $J_0:=\{0\}$) or
$J_{c_\alpha}$, with $c_\alpha\ne c_\omega$. Let $c_1<c_2$ be these two numbers. The orbit of $\bar q$ is disjoint
from $J_{c_1}\cup J_{c_2}$ and accumulates on both, so it lies in the open annulus between them (inside $J_{c_1}$ it
could not accumulate on $J_{c_2}\subset\operatorname{ext}J_{c_1}$, and outside $J_{c_2}$ it could not accumulate on
$J_{c_1}$). Hence $c_1<\bar u<c_2$, and for large $n$ also $c_1<u_n<c_2$, so $J_{u_n}$ separates $J_{c_1}$ from
$J_{c_2}$. The orbit of $\bar q$ is disjoint from the periodic orbit $J_{u_n}$, hence lies inside it or outside it,
and then cannot accumulate on $J_{c_2}$, respectively on $J_{c_1}$. Contradiction; so $\bar u\in\mathrm{Cr}$.
\end{proof}
Let $c_{\rm out}=\max\mathrm{Cr}$, $c_{\rm in}=\min\mathrm{Cr}$, with cycles $\Gamma_{\rm out},\Gamma_{\rm in}$,
$J_{\rm out}=\pi(\Gamma_{\rm out})$, and let $G$ be the bounded component of $\R^2\setminus J_{\rm out}$,
$\overline G=G\cup J_{\rm out}$. Every projected cycle lies in $\overline G$.

\begin{proof}[Proof of (A7)]
\emph{$D\subset\overline G$.} Let $q\in D\setminus\overline G$. Its orbit stays in the exterior of $J_{\rm out}$, so
$\alpha(q)$, $\omega(q)$ lie in $D\setminus G$. By (A6) both are projected cycles ($\{0\}$ is excluded since
$0\in G$); every projected cycle lies in $\overline G$, so both equal $J_{\rm out}$, contradicting (A6).

\emph{$\overline G\subset D$.} Suppose $q_0\in G\setminus D$. Let $y_{\min}=\min_\cA y>0$ and
$g_0=\max(H,y_{\min})$ on $\overline G$; then $g_0=h$ on $J_{\rm out}$, and
$S_0=\{(q,g_0(q)):q\in\overline G\}$ is a compact subset of $\cD$ ($y>0$; and $p<\lambda$ on $\overline G$ because
the linear function $v$ attains its maximum over $\overline G$ on $J_{\rm out}\subset D$). For $t\ge0$ define the
continuous map $\Theta_t:\overline G\to\R^2$, $\Theta_t(q)=\pi\varphi_t(q,g_0(q))$. On $J_{\rm out}$,
$\Theta_t=\psi_t$ maps $J_{\rm out}$ onto itself. Fix $t$ and suppose $q_0\notin\Theta_t(\overline G)$. By
Schoenflies there is a homeomorphism $\chi$ of the closed unit disc onto $\overline G$. The loops
$\Theta_s\circ\chi|_{S^1}$, $0\le s\le t$, lie in $J_{\rm out}\subset\R^2\setminus\{q_0\}$ and form a homotopy from
$\chi|_{S^1}$, which has winding number $\pm1$ about $q_0$ (I4), to $\Theta_t\circ\chi|_{S^1}$; but the latter
extends to the map $\Theta_t\circ\chi$ of the disc into $\R^2\setminus\{q_0\}$ and so has winding number $0$.
Contradiction. Hence for every $n\in\mathbb N$ there is $z_n\in\varphi_n(S_0)$ with $\pi(z_n)=q_0$. By (A1),
$\operatorname{dist}(z_n,\cA)\to0$; a limit point $z\in\cA$ has $\pi(z)=q_0$, so $q_0\in D$. Contradiction.

Thus $D=\overline G$, $\partial D=J_{\rm out}$, and $\cA=(\pi|_\cA)^{-1}(\overline G)$ is a closed topological disc
with boundary $\Gamma_{\rm out}$. The crest statement is Lemma~\ref{lem:crests}.
\end{proof}

\section{Multipliers and corollaries: proof of (A8)}\label{sec:mult}
Let $\Gamma$ be a cycle, $a_0\in\Gamma$, $T$ its period and $\mathsf M$ the monodromy matrix at $a_0$. By
\eqref{eq:normal}, $\Psi(t)=e^{\int_0^t(1+\eta)}n(t)$ solves the adjoint equation $\Psi'=-J^{\mathsf T}\Psi$ along
$\Gamma$, so $\Psi\cdot\delta$ is constant for every solution of $\delta'=J\delta$; with $n$ periodic this gives
$\mathsf M^{\mathsf T}n(a_0)=m_s\,n(a_0)$, $m_s:=e^{-(1+\avg\eta)T}<1$. So $m_s$ is an eigenvalue of $\mathsf M$; another
is $1$ (eigenvector $f(a_0)$), and by Liouville's formula the product of the three is $e^{-AT}$ ($\operatorname{tr}J=-A$).
Hence the third is $m_t=e^{-AT}/m_s=e^{(\avg\eta-\mu)T}$ (this is v2.0, Proposition 2.1). Since
$\eta>\max(0,\mu-1)$ by \eqref{eq:bounds}, $2\avg\eta>\mu-1$, i.e.\ $m_s<m_t$; also $m_s<1$. So $m_s$ is simple and
\begin{equation}\label{eq:nperp}
n(a_0)^\perp=\text{generalised eigenspace of }\mathsf M\text{ for }\{1,m_t\}\ni f(a_0),\qquad n(a_0)\cdot e_s\ne0,
\end{equation}
where $e_s$ is an eigenvector for $m_s$.

\begin{proof}[Proof of (A8)]
(a) By (A5), $\pi(\Gamma)$ is a periodic orbit of $\psi$ with period $T$. If it lies in $\inter D$, $F$ is $C^1$
near it and the classical formula gives the multiplier $\exp\int_0^T\diver F\,dt=\exp\int_0^T(\eta-\mu)\,dt=m_t$
($\eta$ on $\Gamma$ being the function of (A2)). For $\Gamma_{\rm out}$: by (A4) and (I5) there is
$\widehat h\in C^1(\R^2)$ with $\widehat h=h$, $\nabla\widehat h=g$ on $D$. Then $\widehat F=F$ on $D$, each
$\psi$-orbit in $D$ is a $\widehat F$-orbit (uniqueness for the $C^1$ field $\widehat F$), and
$\diver\widehat F=\eta-\mu$ on $D$.

(b) Let $m_t>1$ and suppose $b\in\cA\setminus\Gamma$ has $\omega(b)=\Gamma$. Take a small disc $\Sigma_3$ transverse
to $f$ at $a_0$ with first-return map $P_3$; $DP_3(a_0)$ has eigenvalues $m_s<1<m_t$, with stable eigenvector
$\bar e_s=e_s+c\,f(a_0)$. The orbit of $b$ meets $\Sigma_3$ in points $b_j=P_3^j(b_0)\to a_0$, $b_j\ne a_0$, which by
(I3) lie on the local stable manifold of $a_0$ for $P_3$, a $C^1$ curve tangent to $\bar e_s$. Hence
$(b_j-a_0)/|b_j-a_0|\to\pm\bar e_s/|\bar e_s|$. But $b_j,a_0\in\cA$, so by \eqref{eq:whitney}
\[
\bigl|n(a_0)\cdot(b_j-a_0)\bigr|=h(\pi b_j)-h(\pi a_0)-g(\pi a_0)\cdot(\pi b_j-\pi a_0)\le\omega_g(|b_j-a_0|)\,|b_j-a_0| ,
\]
so $n(a_0)\cdot\bar e_s=0$. By \eqref{eq:nperp} and Lemma~\ref{lem:gauges}(i), $n(a_0)\cdot\bar e_s=n(a_0)\cdot
e_s\ne0$. Contradiction. If $m_t<1$, $\Gamma$ is Lyapunov stable (I3): every neighbourhood $V$ of $\Gamma$ contains a
neighbourhood $V'$ whose forward orbits stay in $V$. If $\alpha(b)=\Gamma$ then $\varphi_{-t_j}(b)\in V'$ for some
$t_j\to\infty$, so $b\in V$ for every $V$, i.e.\ $b\in\Gamma$.

(c) Let $q\in N_E\setminus\{0\}$, so $\alpha(q)=\{0\}$ by (A6). The orbit of $q$ is not periodic, hence disjoint from
$J_{\rm in}$, and it has the point $0\in\operatorname{int}J_{\rm in}$ in its closure; so it lies inside
$J_{\rm in}$. By (A6), $\omega(q)$ is a projected cycle contained in $\overline{\operatorname{int}J_{\rm in}}$; its
crest is $\le c_{\rm in}$, so $\omega(q)=J_{\rm in}$. Thus an orbit in $\cA\setminus\Gamma_{\rm in}$ has
$\omega$-limit set $\Gamma_{\rm in}$, and (b) gives $m_t(\Gamma_{\rm in})\le1$. Suppose
$m_t(\Gamma_{\rm out})>1$. Then $\Gamma_{\rm out}$ is hyperbolic with a two-dimensional local unstable manifold
$W^u_{\rm loc}$, a $C^1$ surface containing $\Gamma_{\rm out}$ whose tangent plane at $a_0\in\Gamma_{\rm out}$ is
spanned by $f(a_0)$ and the eigenvector for $m_t$, i.e.\ equals $n(a_0)^\perp$ by \eqref{eq:nperp}; since $n$ has
third component $-1$, $d\pi$ is injective on it and $\pi(W^u_{\rm loc})$ contains a neighbourhood of
$J_{\rm out}$ in $\R^2$. Points of $W^u_{\rm loc}$ have bounded backward orbits and (by Lemma~\ref{lem:diss})
bounded forward orbits, so $W^u_{\rm loc}\subset\cA$ and $D$ contains points outside $J_{\rm out}$, contradicting
(A7).

(d) Suppose every cycle has $m_t<1$ and there are two cycles. Each crest is isolated in $\mathrm{Cr}$: cycles with
crests $u_n\to c$ have crest points $(u_n,0,h(u_n,0))\to(c,0,h(c,0))$ which are fixed points of the first-return map
to $\{p=0\}$ near the crest point of $J_c$ (each cycle has one crest), and a hyperbolic fixed point is isolated. A
compact set of isolated points is finite, so there are adjacent crests $c_1<c_2$ with no crest between. The open
annulus $R$ between $J_{c_1}$ and $J_{c_2}$ lies in $D$ by (A7) and contains no periodic point. For $q\in R$, (A6)
gives $\alpha(q)\in\{J_{c_1},J_{c_2}\}$ (the orbit stays in $R$, so $\{0\}$ and other cycles are excluded). This
contradicts (b), since both cycles have $m_t<1$.

(e) If all cycles are hyperbolic, $\mathrm{Cr}=\{c_1<\dots<c_N\}$ is finite as in (d). For $q$ in the annulus
between $J_{c_i}$ and $J_{c_{i+1}}$, $\{\alpha(q),\omega(q)\}=\{J_{c_i},J_{c_{i+1}}\}$ by (A6); by (b) the
$\alpha$-limit cycle has $m_t>1$ and the $\omega$-limit cycle has $m_t<1$. So consecutive cycles have opposite
stability, and by (c) the first and last have $m_t<1$; hence $N$ is odd.

(f) On a cycle $\int_0^T\mathbb D\Lambda\,dt=0$, so $\lambda_cT=\int_0^T(\mathbb D\Lambda+\eta-\mu)\,dt<0$ for every
cycle, and (d) applies. With the reverse inequality every cycle would have $\lambda_c>0$, contradicting (c), since
cycles exist.
\end{proof}

\begin{remark}[What is and is not proved]
Theorem~\ref{thm:A} reduces the uniqueness conjecture to the planar system $u'=v$, $v'=k-\mu v-h(u,v)$ on the
closed Jordan domain $D$, with $h$ convex, $h_u=\zeta>0$, $h_v=-\eta<-\max(0,\mu-1)$, the first-order equation
\eqref{eq:pde} and divergence $\eta-\mu$; and it shows that uniqueness at $(\mu,k)$ is \emph{equivalent} to
$\lambda_c<0$ on every cycle there, up to the degenerate case $\lambda_c=0$ (by (A8)(d),(c): uniqueness follows
from $\lambda_c<0$ on all cycles; conversely a unique cycle has $\lambda_c\le0$). It does not decide the sign of
$\avg\eta-\mu$.
\end{remark}

\section{Two rigorous by-products for the uniqueness problem}\label{sec:byproducts}
These statements are proved, but they are \emph{conditional} or \emph{reformulations}: they do not decide uniqueness.

\begin{proposition}[Crest return map]\label{prop:return}
Assume $k>k_H$. Let $\tilde g(u)=h(u,0)-k=k(e^{u}-1)/(1+\eta(u,0))$ and $\tilde G(u)=\int_0^u\tilde g$.
\begin{enumerate}[label=\textup{(\roman*)}]
\item For $0<u<c_{\rm out}$ the $\psi$-orbit of $(u,0)$ returns to $\{v=0,u>0\}$; the first return $R_1(u)$ and the
return time $\tau(u)$ are $C^1$, $R_1$ is increasing, its fixed points are the crests in $(0,c_{\rm out})$, and $R_1$
extends continuously by $R_1(c_{\rm out})=c_{\rm out}$.
\item With $\iota(u)=\int_0^{\tau(u)}(\eta-\mu)(\psi_t(u,0))\,dt$,
\[
R_1'(u)=\frac{\tilde g(u)}{\tilde g(R_1(u))}\,e^{\iota(u)};\qquad\text{in the energy coordinate }e=\tilde G(u):\quad
\frac{d}{de}\,\tilde G\bigl(R_1(\tilde G^{-1}(e))\bigr)=e^{\iota(u)}.
\]
\item If $u\mapsto R_1(u)/u$ is strictly decreasing on $(0,c_{\rm out}]$, or if $\iota$ is strictly decreasing on
$(0,c_{\rm out})$ (equivalently, the energy return map is strictly concave), then the cycle is unique.
\end{enumerate}
\end{proposition}
\begin{proof}
(i) By (A7), $(u,0)\in\inter D\setminus\{0\}$; by (A6) its $\omega$-limit set is a projected cycle $J_c$, and the
$u$-axis is transverse to $F$ at $(c,0)$, so the orbit meets $\{v=0,u>0\}$ at arbitrarily large times. All crossings
of this half-axis are transverse (Remark~\ref{rem:convexgauges}) and $F\in C^1(\inter D)$, so the first return is
$C^1$ by the implicit function theorem; it is increasing because orbits do not cross. Fixed points are periodic
points on the half-axis, i.e.\ crests. Continuity at $c_{\rm out}$ follows from continuity of $\psi$ on $D$ and
transversality at $(c_{\rm out},0)$.
(ii) Let $q(t)=\psi_t(u,0)$ and $\mathsf J(t)=\partial_uq(t)$. The wedge $w=F(q)\wedge\mathsf J$ satisfies
$w'=(\diver F)\,w$; $w(0)=(0,-\tilde g(u))\wedge(1,0)=\tilde g(u)$. Differentiating $\psi_{\tau(u)}(u,0)=(R_1(u),0)$
gives $\mathsf J(\tau)+\tau'(u)F=(R_1'(u),0)$, and wedging with $F(q(\tau))=(0,-\tilde g(R_1(u)))$ gives
$w(\tau)=\tilde g(R_1(u))R_1'(u)$. The second formula follows from $\tilde G'=\tilde g$.
(iii) If $R_1(u)/u$ is strictly decreasing on $(0,c_{\rm out}]$, it equals $1$ only at $c_{\rm out}$. If $\iota$ is
strictly decreasing, the energy map $\tilde R$ is strictly concave on $(0,\tilde G(c_{\rm out}))$ with
$\tilde R(0^+)=0$ (near the equilibrium $0\in\inter D$ the $C^1$ field $F$ has an unstable focus, so return times
stay bounded and $R_1(u)\to0$ as $u\to0^+$), so $\tilde R(e)/e$ is strictly decreasing there and tends to $1$ at the right end point; hence
there is no fixed point in the open interval. In both cases $\Gamma_{\rm out}$ is the only cycle.
\end{proof}
\noindent Numerically (\texttt{job2\_a.py}, \texttt{job2\_a3.py}; 14 parameter pairs) $\iota$ is strictly decreasing
and $R_1$ is concave and sublinear in the coordinate $u$; neither is proved. The analogous turn integral on the
\emph{trough} half-axis is \emph{not} monotone, and $R_1$ is \emph{not} concave in the coordinates $x-k$ or $y-k$.

\begin{remark}[Divergence in other planes of the attractor]\label{rem:planes}
On $\inter D$ the maps $(U,p)\mapsto(U,y)$, $(p,y)$, $(y,m)$, $(U,s)$, with $y=h(U,p)$, $m=p+\mu U$,
$s=p-(y-k)+\mu U$, are local $C^1$ diffeomorphisms (Jacobian determinants, up to sign, $\eta$, $\zeta$, $\zeta+\mu\eta$, $1+\eta$). In
these planes the reduced field is, respectively,
\[
\begin{aligned}
&U'=\mathfrak p(U,y),\ y'=ke^{U}-y;&&\diver=\zeta/\eta-1;\\
&p'=k-\mu p-y,\ y'=ke^{\mathfrak U(p,y)}-y;&&\diver=x/\zeta-A;\\
&y'=ke^{\mathcal U(y,m)}-y,\ m'=k-y;&&\diver=x/(\zeta+\mu\eta)-1;\\
&U'=P(U,s),\ s'=-k(e^{U}-1);&&\diver=(\zeta-\mu)/(1+\eta),
\end{aligned}
\]
where, by convexity of $H$ and the sign of its gradient, $\mathfrak p$ and $P$ are (restrictions of) convex functions
and $\mathfrak U$, $\mathcal U$ of concave functions. Each divergence differs from $\eta-\mu$ by the flow derivative
of the logarithm of the Jacobian: $\mathbb D\log\eta=\zeta/\eta+\mu-1-\eta$, $\mathbb D\log\zeta=x/\zeta-1-\eta$,
$\mathbb D\log(\zeta+\mu\eta)=x/(\zeta+\mu\eta)-(1+\eta-\mu)$, $\mathbb D\log(1+\eta)=(\zeta-\mu)/(1+\eta)+\mu-\eta$.
Hence on every cycle $\lambda_c=\avg\eta-\mu=\avg{\zeta/\eta}-1=\avg{x/\zeta}-A=\avg{x/(\zeta+\mu\eta)}-1
=\avg{(\zeta-\mu)/(1+\eta)}$.
\end{remark}

\section{Numerical consistency checks (not part of the proof)}\label{sec:numerics}
All in ordinary floating point (DOP853, tolerances $10^{-11}$--$10^{-12}$), scripts in the lane directory.
\begin{itemize}
\item \texttt{job1\_checks.py}, parameters $(\mu,\rho)\in\{(1,1.1),(1,2),(0.3,1.5),(3,1.3)\}$, $\rho=k/k_H$.
Pull-back uniqueness: Riccati solutions started from $(10^{-3},10^{-3})$ and $(60,300)$ agree with the reference to
$\le5\times10^{-9}$ after $25$ time units (initial gap $\approx360$). Generalised orthogonality $O=0$ on
$W^u(E)$: relative residual $\le7\times10^{-11}$. Supporting planes \eqref{eq:support} for all ordered pairs among
$3001$ attractor points ($9.0\times10^6$ pairs): no violation in three cases; in the case $(1,1.1)$ one pair of
nearly coincident points gives $\Delta=+3\times10^{-10}$, at the integration-error level. Bounds
\eqref{eq:bounds}: hold. Cone exclusion: no sampled point in $E+C_\pm$; $\min K$, $\min W$,
$\min(y-k-\zeta_EU+\eta_Ep)>0$ on the sample.
\item \texttt{job1\_section.py} (sharp test, event-refined; six parameter pairs). The crest points
$(U,0,y;\eta,\zeta)$ of seven different spirals on $W^u(E)$ are merged and sorted by $U$. In every case $y$ is
strictly increasing in $U$ (graph property) and neighbouring points --- always from \emph{different} spirals ---
satisfy the convexity bracket $\zeta_i\le(y_{i+1}-y_i)/(U_{i+1}-U_i)\le\zeta_{i+1}$, with largest violation
$\le1.2\times10^{-11}$ in $y$-units and none for pairs with $\Delta U>10^{-4}$. The same holds on $\{U=0,p>0\}$
with slope $-\eta$. This checks $\nabla h=(\zeta,-\eta)$ across neighbouring spiral turns.
\item \texttt{job1\_offattractor.py}. A held-out spiral lies on the Hermite-interpolated section curves to within
$5\times10^{-13}k$--$4\times10^{-7}k$. Orbits started at an attractor crest with $y$ shifted by $\delta y$ have
equal $(U,p)$ and distinct $y$ at $t=0$; at their next four section crossings the gap $y-h(U,p)$ equals
$\delta y\,\exp(-\int_0^t(1+\eta))$ within about $1\%$ for $\delta y=\pm0.01k$, the rate predicted by
\eqref{eq:deltaode}.
\item \texttt{job1\_checks2.py}: local plane fits to $y$ over neighbours from other spirals reproduce
$(\zeta,-\eta)$ with an error that decreases with the stencil radius; root-refined crests of successive turns of
$W^u(E)$ increase monotonically to the cycle crest (consistent with (A6)--(A7) in the one-cycle case). A
polyline-crossing test in the same script was discarded as invalid (chord artefacts).
\end{itemize}

\begin{thebibliography}{9}
\bibitem{B} G.~Birkhoff, Extensions of Jentzsch's theorem, \emph{Trans.\ Amer.\ Math.\ Soc.} 85 (1957) 219--227.
\bibitem{Bu} P.~J.~Bushell, Hilbert's metric and positive contraction mappings in a Banach space, \emph{Arch.\
Rational Mech.\ Anal.} 52 (1973) 330--338.
\bibitem{H} P.~Hartman, \emph{Ordinary Differential Equations}, Wiley 1964 (SIAM Classics 2002), Ch.~VII, IX.
\bibitem{S} E.~Seneta, \emph{Non-negative Matrices and Markov Chains}, 2nd ed., Springer 1981, \S3.1.
\bibitem{Sm} H.~L.~Smith, Global dynamics of the smallest chemical reaction system with Hopf bifurcation,
\emph{J.\ Math.\ Chem.} 50 (2012) 989--995.
\bibitem{BH} B.~Boros, J.~Hofbauer, Some minimal bimolecular mass-action systems with limit cycles,
\emph{Nonlinear Anal.\ RWA} 72 (2023) 103839.
\bibitem{W} H.~Whitney, Analytic extensions of differentiable functions defined in closed sets, \emph{Trans.\ Amer.\
Math.\ Soc.} 36 (1934) 63--89.
\end{thebibliography}
\end{document}
